\scanpage{290}
\mspage{16}
% archive sheet number: 289

\unsure{D.} \uline{Axiome} $\mathrm{Loc}(6)$ \uline{des carrés cocartésiens} (ou de ``Mayer-Vietoris''\deleted{\unsure{)}}) (ou des $W$-cofibrations)
% slanted
\margin{C'est le nom ``\uline{axiome des $W$-cofibrations}'' qui m'apparaît à présent comme le meilleur, \unsure{car c'est|en tout cas} \unsure{ce nom|celui} qui me paraît le mieux saisir \unsure{la signification|ce qui} \unsure{conceptuelle|est en jeu} \inserted{\unsure{de l'axiome|en profondeur}}}

J'en vois plusieurs variantes, et n'ai pas encore l'esprit totalement net sur \deleted{\unsure{leurs}} \deleted{\unsure{dépendan}}\unsure{implications} mutuelles. Je commence par la variante plus compliquée, mais dont on aura besoin en tous cas :

\uline{Loc 6} (Axiome \inserted{$W$-cofibrations, ou encore Loc 5, -- axiome} des homomorphismes de carrés cocartésiens). Considérons \inserted{dans $\mathrm{Cat}$} un diagramme commutatif cubique
\diagram{\begin{tikzcd}[ampersand replacement=\&,row sep=3pt,column sep=small]
X_0 \arrow[rr,"{\unsure{i_2|\varepsilon_2}}"] \arrow[dr,"{\unsure{i_1|\varepsilon_1}}"'] \arrow[dd,"f_0"'] \& \& X_2 \arrow[dr] \arrow[dd,"f_2",crossing over] \& \\
\& X_1 \arrow[rr] \arrow[dd,"f_1"'] \& \& X_3 \arrow[dd,"f_3"]\\
Y_0 \arrow[rr,crossing over] \arrow[dr] \& \& Y_2 \arrow[dr] \& \\
\& Y_1 \arrow[rr] \& \& Y_3 ,
\end{tikzcd}}
% slanted
\margin{Il est bon d'introduire des propriétés du diagramme \deleted{\unsure{Loc}} $1^\circ)\ f_0\in W$ \quad $2^\circ)\ f_1$ et $f_2\in W$ \quad $3^\circ)\ f_3\in W$ et de formuler les parties a) b) c) de Loc 6 ainsi \quad a) $1^\circ+2^\circ\Rightarrow3^\circ$ \quad b) $2^\circ+3^\circ\Rightarrow1^\circ$ \quad c) $1^\circ+3^\circ\Rightarrow2^\circ$}
homomorphisme $f=(f_0,f_1,f_2,f_3)$ d'un diagramme carré $X$ dans un autre $Y$.

On suppose que $X$ provient d'un objet $X_3$ de $\mathrm{Cat}$, \deleted{\unsure{muni de deux}}\unsure{de} deux ouverts $X_1$, $X_2$ de $X_3$, \inserted{\uline{recouvrant} $X_3$} et de leur intersection $X_0=X_1\cap X_2$, et de même pour $Y$. (Appelons \inserted{``\uline{couvrants}''} ces diagrammes dans $\mathrm{Cat}$, diagrammes $\mathrm{MV}$ -- ``\uline{Mayer-Vietoris}'')
% slanted
\margin{\uline{NB} On verra plus tard (3.3) que $\uline{\mathrm{Loc}}(6c)$ implique déjà $\mathrm{Loc}(6a,b)$ i.e. $\mathrm{Loc}(6)$ tout entier. Et si $W$ satisfait $\mathrm{Loc}(1\,\mathrm{bis})$ p.~1. \unsure{plus fort|plus fortiori}, alors $\unsure{\mathrm{Loc}(6_0c)}\unsure{\Longleftrightarrow|\Leftarrow|\Rightarrow}\mathrm{Loc}(6a,b)$.}

On \unsure{demande} ce qui suit :
\[
\begin{aligned}
\text{a)}\quad f_0,f_1,f_2\in W &\Longrightarrow f_3\in W\\
\text{b)}\quad f_1,f_2,f_3\in W &\Longrightarrow f_0\in W\\
\bigl[\text{c)}\quad f_0,f_3\in W &\Longrightarrow f_1,f_2\in W\bigr]
\end{aligned}
\]
\margin{NB Il semble que dans $\mathrm{Loc}\,6a$, il suffit de l'\unsure{énoncer|axiome} en supposant $f_0,f_2$ des \uline{iso}, cf p.~22}
\margin{\unsure{Je ne faudrait|Il ne faudrait} pas l'inclure dans $\mathrm{Loc}\,6$ \quad cf p.~25}
